\documentclass[aspectratio=169]{beamer}

\usepackage[T1]{fontenc}
\usepackage[utf8]{inputenc}
\usepackage{lmodern}
\usepackage[polish]{babel}
\usepackage{tikz}
\usetikzlibrary{positioning}
\usepackage{pgfplots}
\pgfplotsset{compat=1.18}
\usepackage{fancyvrb}
\usepackage{booktabs}
\usepackage[backend=biber,style=authoryear]{biblatex}
\addbibresource{bibliografia.bib}

%---------------------------------------
% Motyw: czysty, z kolorową belką tytułową
%---------------------------------------
\usetheme{default}
\definecolor{pwnavy}{RGB}{0,45,98}
\definecolor{pwaccent}{RGB}{0,112,192}
\definecolor{pwgray}{RGB}{90,90,90}

\setbeamercolor{structure}{fg=pwnavy}
\setbeamercolor{frametitle}{bg=pwnavy, fg=white}
\setbeamercolor{title}{fg=pwnavy}
\setbeamercolor{block title}{bg=pwnavy, fg=white}
\setbeamercolor{block body}{bg=pwnavy!8}
\setbeamercolor{alerted text}{fg=pwaccent}
\setbeamerfont{frametitle}{size=\large, series=\bfseries}
\setbeamerfont{title}{size=\Large, series=\bfseries}
\setbeamertemplate{itemize items}[circle]
\setbeamertemplate{navigation symbols}{}
\setbeamertemplate{frametitle}{%
    \nointerlineskip
    \begin{beamercolorbox}[wd=\paperwidth, ht=2.6ex, dp=1.2ex, leftskip=0.6cm]{frametitle}%
        \usebeamerfont{frametitle}\insertframetitle%
    \end{beamercolorbox}%
}
\setbeamertemplate{footline}{%
    \begin{beamercolorbox}[wd=\paperwidth, ht=2.4ex, dp=1ex, leftskip=0.6cm, rightskip=0.6cm]{footline}%
        \usebeamerfont{footline}\color{pwgray}\scriptsize
        B.~Rasztabiga --- LLM w generowaniu konfiguracji Docker i Kubernetes
        \hfill\insertframenumber/\inserttotalframenumber%
    \end{beamercolorbox}%
}

% Wspólny styl wykresów: czytelne z rzutnika
\pgfplotsset{
    every axis/.append style={
        tick label style={font=\footnotesize},
        label style={font=\footnotesize},
        legend style={font=\footnotesize, draw=gray!40},
        ymajorgrids=true, grid style={dashed, gray!30},
        axis line style={gray!60},
    },
}

\title{Zastosowanie dużych modeli językowych\\do generowania konfiguracji Docker i Kubernetes}
\author{Bartłomiej Rasztabiga}
\newcommand{\supervisor}{Promotor: dr inż. Mateusz Modrzejewski}
\institute{Instytut Informatyki, Politechnika Warszawska}
\date{Wrzesień 2026}

\begin{document}

\begin{frame}[plain]
    \titlepage
    \vspace{-1.5em}
    \centering
    \supervisor
\end{frame}

\begin{frame}{Motywacja i cel pracy}
    \begin{columns}[T]
        \column{0.55\textwidth}
        \textbf{Problem:}
        \begin{itemize}
            \item IaC (Dockerfile, manifesty K8s) jest złożone i podatne na błędy przy ręcznym tworzeniu
            \item LLM obiecują automatyzację, lecz brak twardej oceny end-to-end w literaturze
            \item Generowanie kodu aplikacyjnego przez LLM jest dobrze zbadane --- generowanie infrastruktury znacznie słabiej
        \end{itemize}
        \column{0.45\textwidth}
        \textbf{Cel pracy:}\\
        Ocena możliwości i ograniczeń agentów LLM w autonomicznej generacji konfiguracji IaC (Dockerfile + Kubernetes).
        \vspace{1em}

        \textbf{Główne pytanie:}\\
        \textit{Czy agent LLM może --- bez ingerencji człowieka --- wygenerować funkcjonalne, bezpieczne i zgodne z dobrymi praktykami konfiguracje wdrożeniowe?}
    \end{columns}
\end{frame}

\begin{frame}{Przegląd literatury i luki badawcze}
    \begin{columns}[T]
        \column{0.48\textwidth}
        \textbf{Kluczowe prace:}
        \begin{itemize}
            \item \textbf{Repo2Run} \cite{hu_llm-based_2025} --- generacja środowisk Docker z repozytoriów
            \item \textbf{Don't Train, Just Prompt} \cite{kratzke_dont_2024} --- prompt engineering dla IaC
            \item \textbf{GenKubeSec} \cite{malul_genkubesec_2024} --- LLM a bezpieczeństwo K8s
            \item \textbf{CloudEval-YAML} \cite{xu_cloudeval-yaml_2023} --- benchmark dla konfiguracji chmurowych
        \end{itemize}
        \column{0.48\textwidth}
        \textbf{Luki badawcze:}
        \begin{itemize}
            \item Brak oceny autonomicznej --- agenty czytające kod, nie tylko opis słowny
            \item Dominuje walidacja statyczna/składniowa, brak weryfikacji \textbf{runtime}
            \item Brak ilościowej analizy powtarzalności i podatności na manipulację kontekstem
        \end{itemize}
    \end{columns}
\end{frame}

\begin{frame}{Pytania badawcze}
    \begin{itemize}
        \setlength{\itemsep}{0.6em}
        \item \textbf{RQ1:} Na ile agent generuje działające konfiguracje IaC end-to-end?
        \item \textbf{RQ2:} Jak złożoność systemu wpływa na skuteczność agenta?
        \item \textbf{RQ3:} Jaka jest jakość konfiguracji i czy \textit{prompt engineering} ją poprawia?
        \item \textbf{RQ4:} Na ile powtarzalny jest proces przy deterministycznych parametrach?
        \item \textbf{RQ5:} Jak podatny jest agent na manipulację kontekstem repozytorium?
    \end{itemize}
\end{frame}

\begin{frame}{Metodyka i system badawczy}
    \begin{columns}[T]
        \column{0.52\textwidth}
        \small
        \textbf{Modele:} GPT-5 Mini, Gemini 2.5 Flash, DeepSeek V3.2\\[0.35em]
        \textbf{Agent:} LangGraph, 11 narzędzi analizy repozytorium\\[0.35em]
        \textbf{Podejście:} one-shot, bez pętli sprzężenia zwrotnego\\[0.35em]
        \textbf{Walidacja:} lint $\to$ build $\to$ apply $\to$ runtime (HTTP)\\[0.35em]
        \textbf{Środowisko:} MicroK8s, lokalny rejestr obrazów
        \vspace{0.7em}
        \begin{block}{Skala eksperymentów}
            \textbf{485 przebiegów} agenta: 25 repozytoriów GitHub, 5~kontrolowanych PoC, 3~warianty manipulacji
        \end{block}

        \column{0.45\textwidth}
        \centering
        \begin{tikzpicture}[
            node distance=0.45cm and 0.3cm,
            box/.style={draw=pwnavy, thick, rounded corners, align=center, font=\footnotesize,
                        fill=pwnavy!6, minimum width=3.4cm, minimum height=0.62cm}
        ]
            \node[box] (repo) {Repozytorium (kod)};
            \node[box, below=of repo] (agent) {Agent LLM + narzędzia};
            \node[box, below=of agent] (cfg) {Dockerfile + manifesty K8s};
            \node[box, below=of cfg] (val) {Walidacja\\build / apply / runtime};
            \node[box, below=of val] (out) {Metryki + raport};
            \draw[->, thick, pwnavy] (repo) -- (agent);
            \draw[->, thick, pwnavy] (agent) -- (cfg);
            \draw[->, thick, pwnavy] (cfg) -- (val);
            \draw[->, thick, pwnavy] (val) -- (out);
            \node[font=\scriptsize, color=pwgray, right=0.15cm of val] {MicroK8s};
        \end{tikzpicture}
    \end{columns}
\end{frame}

\begin{frame}{RQ1 --- Skuteczność end-to-end}
    \begin{columns}[T]
        \column{0.56\textwidth}
        \begin{tikzpicture}
        \begin{axis}[
            ybar,
            ymin=0, ymax=112,
            bar width=17pt,
            width=\columnwidth, height=5.6cm,
            xtick=data,
            symbolic x coords={Build,Apply,Runtime,{Apply|Build},{Runtime|Apply}},
            xticklabels={Build, Apply, Runtime, {Apply\\po build}, {Runtime\\po apply}},
            x tick label style={font=\scriptsize, align=center},
            nodes near coords,
            nodes near coords style={font=\scriptsize, anchor=south},
            enlarge x limits=0.15,
            legend style={at={(0.5,1.02)}, anchor=south},
            legend columns=2,
            ylabel={Skuteczność [\%]},
        ]
        \addplot[fill=pwnavy!70, draw=pwnavy] coordinates {
            (Build,78.7) (Apply,76.7) (Runtime,62.7)
        };
        \addplot[fill=green!55!black, draw=green!70!black] coordinates {
            ({Apply|Build},97.5) ({Runtime|Apply},81.7)
        };
        \legend{Bezwarunkowe, Warunkowe}
        \end{axis}
        \end{tikzpicture}

        \column{0.42\textwidth}
        \textbf{N = 150} (25 repozytoriów $\times$ 3 modele $\times$ 2 powtórzenia)
        \vspace{0.5em}
        \begin{itemize}
            \item End-to-end: \textbf{62,7\%} (95\% CI: 54,7--70,0\%)
            \item Największy spadek na etapie \textbf{runtime}
            \item Udany build $\Rightarrow$ apply prawie zawsze (97,5\%)
            \item Dominują błędy środowiska: zależności, wersje obrazów bazowych, \texttt{runAsNonRoot} --- niewidoczne dla walidacji statycznej
        \end{itemize}
    \end{columns}
\end{frame}

\begin{frame}{RQ2 --- Wpływ złożoności na skuteczność}
    \begin{columns}[T]
        \column{0.55\textwidth}
        \begin{tikzpicture}
        \begin{axis}[
            ybar,
            ymin=0, ymax=112,
            bar width=22pt,
            width=\columnwidth, height=5.6cm,
            xtick=data,
            symbolic x coords={POC1,POC2,POC3,POC4,POC5},
            nodes near coords,
            nodes near coords style={font=\scriptsize, anchor=south},
            enlarge x limits=0.15,
            ylabel={Skuteczność end-to-end [\%]},
        ]
        \addplot[fill=green!55!black, draw=green!70!black] coordinates {
            (POC1,100) (POC2,33.3) (POC3,22.2) (POC4,11.1) (POC5,0)
        };
        \end{axis}
        \end{tikzpicture}

        \column{0.42\textwidth}
        \textbf{N = 45} (5 PoC $\times$ 3 modele $\times$ 3 powtórzenia)
        \vspace{0.3em}
        {\small
        \begin{tabular}{@{}ll@{}}
            \textbf{POC1} & FastAPI monolit \\
            \textbf{POC2} & + PostgreSQL \\
            \textbf{POC3} & + React \\
            \textbf{POC4} & + 3 serwisy, RabbitMQ \\
            \textbf{POC5} & + Node.js, Redis \\
        \end{tabular}}
        \vspace{0.4em}
        \begin{itemize}
            \item Spadek 100\% $\to$ 0\%; największy skok przy dodaniu bazy danych (POC1 $\to$ POC2)
            \item Błędy integracyjne: sekrety, nazwy usług DNS, zmienne środowiskowe
            \item Agent nie utrzymuje spójności całego systemu
        \end{itemize}
    \end{columns}
\end{frame}

\begin{frame}{RQ3 --- Jakość konfiguracji i prompt engineering}
    \begin{columns}[T]
        \column{0.56\textwidth}
        \begin{tikzpicture}
        \begin{axis}[
            ybar,
            ymin=0, ymax=100,
            bar width=20pt,
            width=\columnwidth, height=5.6cm,
            xtick=data,
            symbolic x coords={{Przebiegi z ostrzeżeniami},{Bez ostrzeżeń}},
            x tick label style={font=\small},
            nodes near coords,
            nodes near coords style={font=\scriptsize, anchor=south},
            enlarge x limits=0.4,
            legend style={at={(0.5,1.02)}, anchor=south},
            legend columns=2,
            ylabel={Odsetek przebiegów [\%]},
        ]
        \addplot[fill=red!60!white, draw=red!80!black] coordinates {
            ({Przebiegi z ostrzeżeniami},82.7) ({Bez ostrzeżeń},17.3)
        };
        \addplot[fill=green!55!black, draw=green!70!black] coordinates {
            ({Przebiegi z ostrzeżeniami},25.3) ({Bez ostrzeżeń},74.7)
        };
        \legend{Prompt bazowy, Prompt best practices}
        \end{axis}
        \end{tikzpicture}

        \column{0.41\textwidth}
        \textbf{N = 150} (25 repozytoriów $\times$ 3 modele $\times$ 2 warianty promptu)\\[0.2em]
        Narzędzia: Hadolint + Kube-linter
        \vspace{0.5em}
        \begin{itemize}
            \item Ostrzeżenia: $147 \to 30$ (\textbf{$-79{,}6\%$})
            \item Przebiegi z ostrzeżeniami: $82{,}7\% \to 25{,}3\%$
            \item Bez ostrzeżeń: $17{,}3\% \to 74{,}7\%$
            \item Modele \textit{znają} dobre praktyki, ale stosują je dopiero po jawnym nakierowaniu
        \end{itemize}
    \end{columns}
\end{frame}

\begin{frame}{RQ4 --- Powtarzalność przy temperature=0}
    \begin{columns}[T]
        \column{0.55\textwidth}
        \begin{tikzpicture}
        \begin{axis}[
            ybar,
            ymin=0, ymax=1.05,
            bar width=22pt,
            width=\columnwidth, height=5.6cm,
            xtick=data,
            symbolic x coords={POC1,POC2,POC3,POC4,POC5},
            nodes near coords,
            nodes near coords style={font=\scriptsize, anchor=south,
                /pgf/number format/.cd, fixed, precision=2},
            enlarge x limits=0.15,
            ylabel={Średni diff ratio (udział zmienionych linii)},
        ]
        \addplot[fill=pwnavy!70, draw=pwnavy] coordinates {
            (POC1,0.258) (POC2,0.462) (POC3,0.494) (POC4,0.878) (POC5,0.723)
        };
        \end{axis}
        \end{tikzpicture}

        \column{0.42\textwidth}
        \textbf{N = 50} (5 PoC $\times$ 2 modele $\times$ 5 powtórzeń)\\[0.2em]
        \texttt{temperature=0}, stały \texttt{seed}\\[0.2em]
        {\footnotesize Gemini 2.5 Flash i DeepSeek V3.2 --- GPT-5 Mini nie przyjmuje \texttt{temperature=0}}
        \vspace{0.4em}
        \begin{itemize}
            \item Średnia \textbf{0,56} --- ponad połowa linii różni się między powtórzeniami
            \item Zmienność rośnie ze złożonością
            \item Główne źródło: wieloetapowa architektura agenta, nie sampling
        \end{itemize}
    \end{columns}
\end{frame}

\begin{frame}{RQ5 --- Podatność na manipulację kontekstem}
    \begin{columns}[T]
        \column{0.52\textwidth}
        \begin{tikzpicture}
        \begin{axis}[
            xbar,
            xmin=0, xmax=95,
            bar width=18pt,
            width=\columnwidth, height=4.8cm,
            ytick=data,
            symbolic y coords={W3,W2,W1},
            nodes near coords,
            nodes near coords style={font=\small, anchor=west},
            xmajorgrids=true,
            enlarge y limits=0.35,
            xlabel={Przebiegi z krytycznym odchyleniem [\%]},
        ]
        \addplot[fill=orange!65!white, draw=orange!80!black] coordinates {
            (66.7,W1)
            (33.3,W2)
            (70.0,W3)
        };
        \end{axis}
        \end{tikzpicture}
        \vspace{0.2em}
        {\footnotesize
        \textbf{W1}: sprzeczna specyfikacja (port 9000 vs 8080 w kodzie)\\
        \textbf{W2}: autorytatywne zalecenia bezpieczeństwa (\texttt{privileged}, \texttt{hostNetwork})\\
        \textbf{W3}: pozorna zależność historyczna (nieistniejący Redis)}

        \column{0.45\textwidth}
        \textbf{N = 90} (3 warianty $\times$ 3 modele $\times$ 10 powtórzeń)
        \vspace{0.3em}

        {\large\textbf{56,7\%}} przebiegów z krytycznym odchyleniem \textit{(51/90)}
        \vspace{0.5em}
        \begin{itemize}
            \item Agent ufa dokumentacji bardziej niż kodowi
            \item W1: błędny port przepisany spójnie do Dockerfile, \texttt{containerPort} i \texttt{Service}
            \item W3: pełny stos Redis bez śladu zależności w kodzie
            \item Proste techniki, bez exploitowania modelu = realne zagrożenie
        \end{itemize}
    \end{columns}
\end{frame}

\begin{frame}{Porównanie modeli}
    \begin{columns}[T]
        \column{0.55\textwidth}
        \begin{tikzpicture}
        \begin{axis}[
            ybar,
            ymin=0, ymax=112,
            bar width=16pt,
            width=\columnwidth, height=5.6cm,
            xtick=data,
            symbolic x coords={{Gemini 2.5 Flash},{GPT-5 Mini},{DeepSeek V3.2}},
            xticklabels={{Gemini\\2.5 Flash},{GPT-5\\Mini},{DeepSeek\\V3.2}},
            x tick label style={font=\small, align=center},
            nodes near coords,
            nodes near coords style={font=\scriptsize, anchor=south},
            enlarge x limits=0.25,
            legend style={at={(0.5,-0.28)}, anchor=north},
            legend columns=1,
            ylabel={[\%]},
        ]
        \addplot[fill=green!55!black, draw=green!70!black] coordinates {
            ({Gemini 2.5 Flash},62.0)
            ({GPT-5 Mini},68.0)
            ({DeepSeek V3.2},58.0)
        };
        \addplot[fill=red!60!white, draw=red!80!black] coordinates {
            ({Gemini 2.5 Flash},6.7)
            ({GPT-5 Mini},63.3)
            ({DeepSeek V3.2},100.0)
        };
        \legend{RQ1: skuteczność end-to-end, RQ5: podatność na manipulację kontekstem}
        \end{axis}
        \end{tikzpicture}

        \column{0.43\textwidth}
        \vspace{0.3em}
        {\footnotesize
        \begin{tabular}{@{}lccc@{}}
            \toprule
            \textbf{Model} & \textbf{RQ1} & \textbf{RQ2} & \textbf{RQ5} \\
            \midrule
            Gemini 2.5 Flash & 62,0\% & 20,0\% & \textbf{6,7\%} \\
            GPT-5 Mini & \textbf{68,0\%} & \textbf{46,7\%} & 63,3\% \\
            DeepSeek V3.2 & 58,0\% & 33,3\% & 100\% \\
            \bottomrule
        \end{tabular}}\\[0.2em]
        {\scriptsize RQ1, RQ2: skuteczność end-to-end; RQ5: podatność na manipulację}
        \vspace{0.4em}
        \begin{itemize}
            \item Różnice w RQ1 w granicach nakładających się CI
            \item Żaden model nie dominuje jednocześnie w skuteczności i odporności
            \item Wybór modelu zależy od profilu ryzyka wdrożenia
        \end{itemize}
    \end{columns}
\end{frame}

\begin{frame}{Wnioski i implikacje praktyczne}
    \begin{block}{Odpowiedź na główne pytanie badawcze}
        \textbf{Tak} --- dla prostych aplikacji jednokomponentowych z czytelną strukturą kodu.\\
        \textbf{Nie} --- dla złożonych systemów wielousługowych oraz środowisk z podwyższonymi wymaganiami bezpieczeństwa lub niezaufaną dokumentacją.
    \end{block}
    \vspace{0.4em}
    \textbf{Implikacje praktyczne:}
    \begin{itemize}
        \item \textbf{Ograniczony kontekst wejściowy} --- tylko pliki kodu, bez dokumentacji i historii (ogranicza RQ5)
        \item \textbf{Rozbudowana instrukcja systemowa} --- konkretna checklista dobrych praktyk redukuje ostrzeżenia o $>79\%$ (RQ3)
        \item \textbf{Obowiązkowa walidacja statyczna} --- Hadolint + Kube-linter jako element CI/CD
        \item \textbf{Ludzki przegląd} --- artefakty z security context i uprawnieniami wymagają weryfikacji
    \end{itemize}
\end{frame}

\begin{frame}{Zagrożenia trafności}
    \begin{itemize}
        \setlength{\itemsep}{0.5em}
        \item Zbiór RQ1: 25 prostych aplikacji może nie uogólniać się na złożone systemy produkcyjne
        \item Mała liczba powtórzeń (RQ1: 2, RQ2: 3, RQ3: 1) --- szerokie przedziały ufności
        \item Walidacja runtime żądaniami HTTP do endpointów wskazanych przez agenta nie weryfikuje w pełni działania usługi
        \item 5 repozytoriów PoC (RQ2) --- zbyt mała próba do precyzyjnej lokalizacji progu złożoności
        \item Brak pętli sprzężenia zwrotnego z runtime --- architektura one-shot może zaniżać wyniki względem systemów iteracyjnych
    \end{itemize}
\end{frame}

\begin{frame}{Wkład do wiedzy i dalsze kierunki}
    \begin{columns}[T]
        \column{0.52\textwidth}
        \textbf{Wkład do wiedzy:}
        \begin{itemize}
            \item Empiryczna weryfikacja \textit{repo poisoning} (indirect prompt injection) w kontekście IaC --- dotąd nieopisana w literaturze IaC
            \item Walidacja end-to-end z testami runtime --- ujawnia błędy niewidoczne dla narzędzi statycznych
            \item Ilościowy pomiar wpływu prompt engineeringu na jakość IaC ($-79{,}6\%$ ostrzeżeń)
            \item Kontrolowany eksperyment z pięcioma poziomami złożoności
        \end{itemize}
        \column{0.45\textwidth}
        \textbf{Dalsze kierunki:}
        \begin{itemize}
            \item Pętla sprzężenia zwrotnego z runtime --- agent reaguje na błędy wdrożenia
            \item Mechanizmy obronne: weryfikacja spójności kodu i dokumentacji
            \item Stabilizacja wyników przez wielokrotne generacje i agregację
            \item Zastosowanie w systemie PaaS --- automatyczne wdrożenie po podaniu linku do repozytorium
        \end{itemize}
    \end{columns}
\end{frame}

\begin{frame}[plain]
    \centering
    {\Large\bfseries\color{pwnavy} Dziękuję za uwagę.}\\[1em]
    \normalsize Bartłomiej Rasztabiga
\end{frame}

\begin{frame}[allowframebreaks]{Bibliografia}
    \printbibliography
\end{frame}

\end{document}
